% Part: many-valued-logic
% Chapter: syntax-and-semantics
% Section: introduction

\documentclass[../../../include/open-logic-section]{subfiles}

\begin{document}

\olfileid{mvl}{syn}{int}

\olsection{مقدمة}

في المنطق الكلاسيكي نبني~!!{formula}s من~!!{propositional variable}s
بروابط القضايا~$\lnot$ و$\land$ و$\lor$ و$\lif$ و$\liff$.
ثم نجعل لدلالته قيمتي الصدق~$\True$ و$\False$. فـ
  !!{propositional variable}s تفسرها~!!a{valuation}~$\pAssign{{\OLId{latin-ordinary}{v}}}$
  تسند إلى كل واحد منها إحدى القيمتين~$\True$ و$\False$؛ أي لا تسند
  إلى~!!{propositional variable}s إلا قيمة واحدة لكل واحد منها. ومن هذا
التعيين يحدد كل~!!{valuation} قيمة~$\pValue{{\OLId{latin-ordinary}{v}}}(!A)$ لكل
!!{formula}~$!A$. وتكون~!!^a{formula} مشبعة في~!!a{valuation}
$\pAssign{{\OLId{latin-ordinary}{v}}}$، ويرمز إلى ذلك بـ~$\pSat{{\OLId{latin-ordinary}{v}}}{!A}$، إذا وفقط إذا كان
$\pValue{{\OLId{latin-ordinary}{v}}}(!A) = \True$.

وأما المنطقات المتعددة القيم فتعمم هذا البناء الثنائي، إذ تجيز قيم
صدق سوى~$\True$ و$\False$. فيكون~!!a{valuation}~$\pAssign{{\OLId{latin-ordinary}{v}}}$ فيها
دالة تسند إلى كل~!!{propositional variable}~${\OLId{latin-ordinary}{p}}$ عنصرًا من مجموعة
قيم الصدق الجائزة، ولنسم هذه المجموعة عادة~${\OLId{latin-looped}{V}}$. وبهذا الاعتبار يكون
المنطق الكلاسيكي نفسه منطقًا متعدد القيم، ومجموعة قيم صدقه
${\OLId{latin-looped}{V}} = \{\True, \False\}$؛ وتحسب فيه~$\pValue{{\OLId{latin-ordinary}{v}}}(!A)$ بجداول الصدق
المألوفة للروابط.

% تصحيحات أسلوبية (0385-M1، 0385-N1، 0385-N2): صُحح تركيب الإضافة، وفُصلت مجموعة القيم عن المصفوفة، وأزيل إيهام إسناد القيمتين معًا.

فإذا وسعنا قيم الصدق لم يعد تفسير الروابط التي نقرأها «و» و«أو»
و«ليس» و«إذا---فإن» متعينًا على وجه واحد؛ بل ظهرت وجوه طبيعية شتى
لحساب~$\pValue{{\OLId{latin-ordinary}{v}}}(!A)$. ولذلك لا يكفي في تعيين المنطق المتعدد القيم
أن نذكر مجموعة قيمه، بل لا بد من تعيين \emph{دوال الصدق} المختارة
لكل رابط. فإذا اخترناها لمنطق~$\Log L$، عين كل تقييم قيمة
$\pValue{{\OLId{latin-ordinary}{v}}}(!A)[\Log L]$ تعيينًا وحيدًا، كما في المنطق الكلاسيكي.
فالمنطقان، من هذه الجهة، \emph{داليّا الصدق}.

وبعد ضبط هذه الأصول الدلالية نعرّف نظائر التحصيل المنطقي والاستلزام
وقابلية الإشباع. ففي المنطق الكلاسيكي تكون~!!a{formula} تحصيلًا
منطقيًا إذا تحقق~$\pValue{{\OLId{latin-ordinary}{v}}}(!A) = \True$ لكل~$\pAssign{{\OLId{latin-ordinary}{v}}}$.
أما في المنطق المتعدد القيم فلا بد من تعميم آخر؛ إذ قد تخلو مجموعة
قيم الصدق~${\OLId{latin-looped}{V}}$ من القيمة~$\True$ أصلًا. فقد تتخذ بعض المنطقات، مثلًا، جميع الأعداد
النسبية بين~${\OLMathNumeral{0}}$ و${\OLMathNumeral{1}}$ قيمًا للصدق، وحينئذ يقوم~${\OLMathNumeral{1}}$ عادة مقام~$\True$.
وقد تقوم بهذا المقام في منطق آخر عدة قيم، لا قيمة واحدة. فنخص كل
منطق متعدد القيم بمجموعة~${\OLId{latin-looped}{V}}^+$ من \emph{القيم المميزة}. وعندئذ تكون
!!a{formula} مشبعة في~!!a{valuation}~$\pAssign{{\OLId{latin-ordinary}{v}}}$، أي
$\pSat{{\OLId{latin-ordinary}{v}}}{!A}[\Log L]$، إذا وفقط إذا كان
$\pValue{{\OLId{latin-ordinary}{v}}}(!A)[\Log L] \in {\OLId{latin-looped}{V}}^+$. وتكون~!!^a{formula}~$!A$
تحصيلًا منطقيًا لذلك المنطق، أي~$\Entails[\Log L] !A$، إذا وفقط إذا
كان~$\pValue{{\OLId{latin-ordinary}{v}}}(!A)[\Log L] \in {\OLId{latin-looped}{V}}^+$ لكل~$\pAssign{{\OLId{latin-ordinary}{v}}}$. وأخيرًا تكون~$!A$
مستلزمة من مجموعة~!!{formula}s، أي~${\OLId{greek-variable}{Gamma}} \Entails[\Log L] !A$،
إذا كان كل~!!{valuation} يشبع جميع~!!{formula}s في~${\OLId{greek-variable}{Gamma}}$ يشبع
$!A$ أيضًا.

\end{document}
